\section{Introduction}\label{sec:intro}
Multi-robot inspection, communication relaying, and formation reconfiguration require collision avoidance, network connectivity, and task completion. Algebraic connectivity characterizes a state-dependent communication graph, and distributed estimation and control provide tools for maintaining mobile networks\cite{yang2010decentralized,sabattini2013decentralized,ong2021connectivity}. Control barrier functions (CBFs), often combined with quadratic programs (QPs), impose safety conditions while modifying nominal commands\cite{ames2017cbf,breeden2022sampled}. Joint spectral and pairwise constraints additionally require allocation of coupled responsibilities under finite communication\cite{tan2025violationfree}.

Instantaneous feasibility does not imply task completion. CBF-QPs can introduce undesired equilibria, and robots exchanging positions can become stuck\cite{reis2021,tan2024equilibria,grover2023deadlock}. For $W(\bp)=\frac12\sum_i\norm{\bp_i-\bg_i}_2^2$, projecting $-\nabla W$ onto a convex safe input set containing zero produces error-monotone control\cite{allibhoy2024safegradient}. A strict nongoal local minimum on the feasible configuration set can require a temporary increase in $W$ along a successful safe path. Robots already at their own goals may need to move to permit interior rearrangement.

Analytical approaches introduce circulation or compatible control Lyapunov and barrier conditions\cite{goncalves2023circulation,dai2024compatible}. Learning approaches use constrained policy optimization, safety filters, or learned barrier functions\cite{gu2023safe,wabersich2021}. GCBF+ learns neural graph barrier functions and distributed control\cite{zhang2025gcbf}; DGPPO incorporates discrete graph barrier conditions into proximal policy updates\cite{zhang2025dgppo}. Their constraint representations, dynamics, observations, and objectives differ, so empirical comparisons require explicit task adaptations. Static-allocation and custom projection-penalty experiments describe those internal variants, not these independent published methods.

We investigate three connected design questions: generating cooperative nonmonotone motion when error descent is obstructed, recruiting stationary neighbors through local allocation, and suppressing learned residuals at the goal. The proposed Envelope Equivariant Graph Policy with Neighborhood Proximal Allocation (EGP-NPA) combines multilayer neighbor-message passing, a coupled safety-allocation layer, and a completion envelope restoring the zero-velocity target equilibrium.

First, a rational second-order Karush--Kuhn--Tucker (KKT) certificate and a safe path on a seven-robot contact configuration expose a limitation of error-monotone control. Second, we construct a rotation-, translation-, and permutation-equivariant residual policy, analyze deterministic lateral-output restrictions on exactly mirror-symmetric configurations, and derive allocation-induced yielding in a two-node contact subproblem. Third, we state conditional safety guarantees for learned commands followed by allocation and holding-segment certification, and establish a first-order envelope bound. Finally, specified-checkpoint evaluations, fixed-weight interventions, and budget- and seed-matched learning comparisons distinguish controller behavior, component effects, and training variation. Negative second-order physical-simulation outcomes are retained.

\textbf{Notation:} $\mathbb R^n$, $\mathbb Q^n$, and $\mathbb S_+^n$ denote real vectors, rational vectors, and positive semidefinite matrices. $\I_n$ is the identity, $\boldsymbol1_N$ is the all-ones vector, and $\boldsymbol\varepsilon_i$ is a standard basis vector. $\norm{\cdot}_2$ denotes the Euclidean or spectral norm; $\col$ denotes column stacking; $\otimes$ denotes the Kronecker product; $(x)_+=\max(x,0)$ and $\clip(x,a,b)=\min(\max(x,a),b)$. The matrix $\bJ=\left[\begin{smallmatrix}0&-1\\1&0\end{smallmatrix}\right]$ rotates vectors counterclockwise by a right angle, and $\mathbb I(\cdot)$ is an indicator.
